% ==============================================================================
% РАЗДЕЛ 3.7 (DISS-005-B2047 .. DISS-005-B2053)
% Выводы по главе 3
% Замечания: DISS-005-C0136 (1 шт.)
% ==============================================================================

% DISS-005-B2047
\section{Выводы}
\label{sec:ch03_conclusions}

% DISS-005-B2048
В настоящей главе был детально описан комплексный и многоэтапный процесс формирования и подготовки обширного датасета реальных осциллограмм, предназначенного для широкого круга исследований в области релейной защиты, автоматики и анализа переходных процессов в электроэнергетических системах. Были представлены разработанные методики, охватывающие все стадии жизненного цикла данных от сбора из разнообразных источников и первичной обработки, включающей унификацию множества форматов и исключение дубликатов, до глубокой предобработки, такой как деперсонализация для обеспечения конфиденциальности, стандартизация имен сигналов, определение базисных значений для корректной нормализации и фильтрация неинформативных («пустых») записей.

% DISS-005-B2049
Результатом проделанной работы стал структурированный и публично доступный датасет, содержащий 50~765 реальных осциллограмм. Уникальность данного датасета заключается не только в его объеме, но и в разнообразии представленных данных, охватывающих широкий спектр аварийных событий (включая различные виды коротких замыканий и однофазные замыкания на землю), оперативных переключений, специфических режимов работы оборудования (например, пусковые режимы двигателей, работа устройств БАВР), а также шумовых записей. Разработанная методология, включающая автоматизированный поиск, копирование с учетом уникальности (MD5-хеширование), иерархическую классификацию и аннотирование событий, обеспечивает высокое качество, репрезентативность и масштабируемость созданной базы данных.

% DISS-005-B2050
Сформированный и тщательно подготовленный датасет является критически важным фундаментом для дальнейших исследований, представленных в последующих главах настоящей диссертации, а также для научного сообщества в целом.

% DISS-005-B2051
Во-первых, он открывает широкие возможности для проведения глубокого статистического анализа реальных процессов в энергосистемах (что будет рассмотрено в Главе 4\commentnote{DISS-005-C0136}{Алексей Евдаков [2]}{Глава 4 и 5 пока плановые, возможно будут заметко скорректированы.\par\vspace{0.1cm}Пока план такой:\par--~Глава 4. статистическая обработка текущего датасета.\par--~Глава 5. разработка нейронок для разметки данных (очень вероятно перейдёт в главу 4) и разработка современных органов РЗА, а точнее РНМ на основе нейронных сетей.}). На основе этого датасета будет проведен анализ временных характеристик используемых выключателей, оценка реальных перенапряжений при однофазных замыканиях на землю, что позволит верифицировать существующие теоретические модели и выявить новые, ранее не описанные закономерности.

% DISS-005-B2052
Во-вторых, данный датасет, особенно его аннотированная часть, является необходимой и достаточной основой для разработки и обучения моделей машинного обучения (чему будет посвящена Глава 5). Первоначальная ручная разметка ключевых событий и разработанные подходы к автоматизации этого процесса, включая применение нейронных сетей для классификации и сортировки осциллограмм, позволят не только существенно ускорить анализ больших объемов данных, но и перейти к созданию современных адаптивных алгоритмов РЗА. В частности, планируется разработка и исследование новых подходов к реализации реле направления мощности (РНМ) на основе нейронных сетей, способных эффективно функционировать в сложных и быстроменяющихся условиях эксплуатации.

% DISS-005-B2053
Проделанная работа по созданию и подготовке датасета не только решает проблему нехватки качественных реальных данных для исследований, но и закладывает прочную основу для разработки инновационных методов анализа и управления в электроэнергетике. Дальнейшие перспективы развития включают расширение датасета за счет данных из сетей более высокого класса напряжения и сетей с частотой 60~Гц. А также углубление и детализацию аннотаций, и применение методов обучения без учителя для выявления скрытых паттернов и аномалий в поведении энергосистем.
